%% ma: L12-B14-0049
%% lop: 12
%% chuong: 5
%% bai: 14
%% dang: 12.14.1
%% loai: TN4
%% muc: TH
%% dap_an: "A"
%% nguon: "To_n_12_PT_M_t_Ph_ng.pdf (Đề kiểm tra đánh giá kiến thức, Chuyên đề Phương trình mặt phẳng, Nguyễn Anh Khôi - Lê Minh Khôi), Câu 3"
%% kien_thuc: "Bài 14 (phương trình mặt phẳng theo đoạn chắn)"
%% trang_thai: chinh-thuc
%% ghi_chu: "Tài liệu gốc không kèm đáp án; đáp án do Claude giải và kiểm bằng sympy."
\begin{ex}
Trong không gian $Oxyz$, một vectơ pháp tuyến của mặt phẳng $\dfrac{x}{-2}+\dfrac{y}{-1}+\dfrac{z}{3}=1$ là
\choice{\True $\vec n=(3;6;-2)$}{$\vec n=(2;-1;3)$}{$\vec n=(-3;-6;-2)$}{$\vec n=(-2;-1;3)$}
\loigiai{Nhân hai vế với $-6$: $3x+6y-2z+6=0$. Vectơ pháp tuyến $\vec n=(3;6;-2)$. Chọn A.}
\end{ex}
